\chapter{Notation Summary}
\label{ch:app-c-notation}

This appendix gathers the mathematical notation used across the book in one reference, defines each macro once, and states the intended reading beside it, so later chapters can reuse these symbols without redefinition.
A chapter that needs a different meaning for an existing symbol must introduce a distinct macro instead of silently reusing the name.

\section{Dataflow facts and transfer}
\label{sec:app-c-dataflow}

Dataflow notation names facts that hold at block boundaries and the functions that move those facts across blocks.
Table~\ref{tab:app-c-dataflow} lists the macros and their one-line readings.
\begin{table}
  \centering
  \caption{Dataflow macros name boundary facts, generated facts, and lattice operations.}
  \label{tab:app-c-dataflow}
  \begin{tabular}{ll}
    \toprule
    Macro & Meaning \\
    \midrule
    \texttt{\textbackslash In\{B\}} & fact true on entry to block $B$ \\
    \texttt{\textbackslash Out\{B\}} & fact true on exit from block $B$ \\
    \texttt{\textbackslash Gen\{B\}} & fact created inside block $B$ \\
    \texttt{\textbackslash Kill\{B\}} & fact removed by block $B$ \\
    \texttt{\textbackslash LiveIn\{B\}} & variable live on entry to block $B$ \\
    \texttt{\textbackslash LiveOut\{B\}} & variable live on exit from block $B$ \\
    \texttt{\textbackslash Tf\{B\}} & transfer function of block $B$ \\
    \texttt{\textbackslash meet} & greatest lower bound of lattice facts \\
    \texttt{\textbackslash join} & least upper bound of lattice facts \\
    \bottomrule
  \end{tabular}
\end{table}
$\In{B}$ and $\Out{B}$ describe the analysis result at the two edges of a block.
$\Gen{B}$ and $\Kill{B}$ summarize what the block itself contributes before any fixed-point iteration runs.
$\LiveIn{B}$ and $\LiveOut{B}$ specialize the same idea to the liveness question of whether a variable may be read later.
$\Tf{B}$ maps entry facts to exit facts in a forward analysis, and exit facts to entry facts in a backward analysis.
The symbols $\meet$ and $\join$ combine facts where control-flow paths meet.
Their exact direction depends on the analysis, so a chapter must state which operator it uses and why.
Chapters~\ref{ch:ch10-lattices-fixed-points} and~\ref{ch:ch14-constant-propagation} use a reversed precision order for flat constants: top is the unvisited identity and bottom admits conflicting values.
Chapter~\ref{ch:ch11-abstract-interpretation} uses the usual concretization order: bottom is unreachable and top admits every value.
These conventions are dual, not interchangeable; translate the order, extrema, and merge operator together when moving an analysis between them.
All of these macros keep the same typographic form in every chapter, so a reader can carry that reading into later dataflow clients.
Figure~\ref{fig:visual-subset-lattice} anchors the symbols to a concrete four-element example.
Its arrows mean subset inclusion; reversing the order would exchange which operation is meet and which is join.
The diagram therefore belongs with its stated order, not just with the names of the four sets.

\begin{figure}[htbp]
  \centering
  \begin{tikzpicture}[node distance=12mm and 16mm]
    \node[basic block] (top) {$\{p,q\}$};
    \node[basic block,below left=of top] (p) {$\{p\}$};
    \node[basic block,below right=of top] (q) {$\{q\}$};
    \node[basic block,below right=of p] (bot) {$\emptyset$};
    \draw[cfg edge] (bot) -- (p);
    \draw[cfg edge] (bot) -- (q);
    \draw[cfg edge] (p) -- (top);
    \draw[cfg edge] (q) -- (top);
    \node[font=\small,anchor=west,right=12mm of q]
      {$\{p\}\join\{q\}=\{p,q\}$};
    \node[font=\small,anchor=west,right=12mm of bot]
      {$\{p\}\meet\{q\}=\emptyset$};
  \end{tikzpicture}
  \caption{For subsets ordered by inclusion, upward arrows add possibilities, join is union, and meet is intersection.}
  \label{fig:visual-subset-lattice}
\end{figure}



\section{SSA, control flow, and definitions}
\label{sec:app-c-ssa-cfg}

Static single assignment and control-flow notation names where values are defined and how blocks connect.
Table~\ref{tab:app-c-ssa-cfg} lists the macros and their one-line readings.
\begin{table}
  \centering
  \caption{SSA and control-flow macros name merges, predecessors, dominance, and definitions.}
  \label{tab:app-c-ssa-cfg}
  \begin{tabular}{ll}
    \toprule
    Macro & Meaning \\
    \midrule
    \texttt{\textbackslash phifn\{...\}} & merge of values from predecessor blocks \\
    \texttt{\textbackslash assignphi} & assignment of a merged value to a name \\
    \texttt{\textbackslash Pred} & set of CFG predecessors of a block \\
    \texttt{\textbackslash Succ} & set of CFG successors of a block \\
    \texttt{\textbackslash dom} & set of blocks dominating a block \\
    \texttt{\textbackslash idom} & immediate dominator of a block \\
    \texttt{\textbackslash DF} & dominance frontier of a block \\
    \texttt{\textbackslash Def} & set of definitions of a variable \\
    \texttt{\textbackslash Use} & set of uses of a variable \\
    \texttt{\textbackslash Entry} & distinguished CFG entry block \\
    \texttt{\textbackslash Exit} & distinguished CFG exit block \\
    \bottomrule
  \end{tabular}
\end{table}
$\phifn{x_1,x_2}$ selects among incoming values according to the edge actually taken.
$\Pred$ and $\Succ$ describe graph adjacency without implying anything about execution frequency.
The relations $\dom$, $\idom$, and $\DF$ describe which blocks necessarily precede or rejoin the flow through a given block.
$\Def$ and $\Use$ connect a variable occurrence to the program points that write or read it.
$\Entry$ and $\Exit$ anchor analyses that propagate facts forward or backward across the graph.
These macros appear in figure labels, algorithm statements, and equations with the same shapes every time, so a reader can recognize a dominance claim or a merge site at a glance.

\section{Sets and evaluation in the opening chapter}
\label{sec:app-c-sets-traces}

Set-builder and tiny-language notation appear heavily in the opening chapter and recur whenever a later chapter reasons about values.
Table~\ref{tab:app-c-ch01} lists the symbols and their one-line readings.
\begin{table}
  \centering
  \caption{Set, evaluation, translation, and behavior symbols used in the opening chapter.}
  \label{tab:app-c-ch01}
  \begin{tabular}{ll}
    \toprule
    Symbol & Meaning \\
    \midrule
    \texttt{\textbackslash set\{x | P(x)\}} & set of values $x$ satisfying property $P$ \\
    $V$ & tiny-language value set $\{0,1,\ldots,255\}$ \\
    $\llbracket e\rrbracket_a$ & value of expression $e$ at input $a$ \\
    $F(e)$ & constant-folded form of expression $e$ \\
    $C(e)$ & stack instruction sequence compiled from $e$ \\
    $\operatorname{Exec}(c,a,S)$ & result of running code $c$ at input $a$ from stack $S$ \\
    $S\mathbin{\|} [v]$ & stack $S$ extended with value $v$ \\
    $D(e)$ & peak extra stack slots needed by $C(e)$ \\
    $(t,r)$ & terminating execution with trace $t$ and result $r$ \\
    $\mathcal{B}(P,i)$ & set of behaviors permitted for program $P$ on input $i$ \\
    $\operatorname{run}(C(P),i)$ & observed execution of compiled program on input $i$ \\
    \bottomrule
  \end{tabular}
\end{table}
The macro $\set{x \in S | x > 0}$ builds a set from a membership condition, with the bar separating the introduced variable from its predicate.
The value set $V$ fixes the pure tiny-language domain; later languages state their own value domains.
The bracket form $\llbracket e\rrbracket_a$ means tree evaluation under the semantics in force, including trap propagation in the division extension.
The letter $F$ always means the constant folder and the letter $C$ always means the stack-code compiler in that chapter.
The executor form $\operatorname{Exec}(c,a,S)$ makes the initial stack explicit so fragments and whole programs can be distinguished.
The concatenation symbol joins either stacks or instruction sequences, with the surrounding expression identifying which kind is meant.
The demand function $D(e)$ computes the exact peak from the tree shape without running the machine.
The pair $(t,r)$ separates the observable event sequence from the final value.
The behavior set $\mathcal{B}(P,i)$ collects every execution the language permits for one program and input.
The refinement requirement $\mathcal{B}(C(P),i)\subseteq\mathcal{B}(P,i)$ then says that compilation must not introduce a forbidden behavior.
The simpler equality $\operatorname{run}(C(P),i)=\operatorname{run}(P,i)$ applies only when both executions have plain results and those results capture everything worth preserving.

\section{Reading conventions}
\label{sec:app-c-conventions}

Mathematics always appears in math mode, including single variables such as $x$, $a$, $u$, and $v$.
Code identifiers, opcodes, and file names use \code{...} rather than math italics.
The distinction keeps the tiny-language variable $x$ visually separate from a host variable named \code{x} where that difference matters.
Block names such as $B$, $P$, and $Q$ are metavariables ranging over blocks or expressions.
Subscripts select an instance, as in $\In{B}$ for the entry fact of block $B$.
The symbols $\meet$ and $\join$ denote lattice operations; a particular set lattice may implement one of them as union or intersection.
Parentheses in $\phifn{x_1,x_2}$ list merge inputs in predecessor order.
Square brackets in $S\mathbin{\|}[v]$ delimit stack contents rather than indexing an array.
Mathematical equality states that two expressions denote the same value.
The refinement inclusion states that one behavior set is contained in another.
Conflating those two relations would turn a permission argument into an equality claim.
A \term{conservative analysis} keeps enough possibilities to avoid an unsound claim, and its direction depends on what the analysis represents.
For call targets, retaining an extra target stays safe while omitting a real target does not.
For constants, claiming a value without justification is unsound while keeping the original computation is merely slower.
This appendix remains the canonical place to check the spelling and shape of every shared symbol.
